\chapter{Shortness of Breath}

\section*{Definition}
Shortness of breath is the uncomfortable feeling of struggling to breathe or not getting enough air. It can range from mild to severe.

\section*{When to See a Doctor}
Seek emergency care immediately for severe or sudden breathlessness. Otherwise arrange prompt medical assessment if symptoms are new, recurrent, or worsening:
\begin{itemize}
  \item Sudden shortness of breath, blue or grey lips, confusion, fainting, inability to speak full sentences, or severe chest pain
  \item Breathlessness at rest, rapidly worsening symptoms, or markedly low oxygen readings if a reliable pulse oximeter is available
\end{itemize}

\section*{How It Develops}
The sensation of shortness of breath arises from a mismatch between what your brain expects and what your body delivers. The cause may involve your heart, lungs, blood, metabolism, or even anxiety.

\section*{Causes}
\begin{itemize}
  \item Heart failure exertion shortness of breath when lying flat
  \item Asthma COPD exacerbation
  \item Pulmonary embolism acute onset
  \item Pneumonia infection
  \item Anxiety panic attack hyperventilation
\end{itemize}

\section*{Related Symptoms}
\begin{itemize}
  \item Chest tightness wheezing cough shortness of breath when lying flat paroxysmal nocturnal shortness of breath
\end{itemize}

\section*{Medical Evaluation}
\begin{itemize}
  \item History focuses on onset, duration, character of symptoms (productive vs dry cough, nasal discharge color), fever pattern, and exposure history.
  \item Physical examination includes vital signs, lung auscultation, nasal inspection, and oropharyngeal examination.
  \item Targeted testing may include rapid antigen/ PCR for respiratory viruses, chest X-ray, pulse oximetry, and sputum culture.
  \item Allergy testing or pulmonary function tests are reserved for chronic or recurrent cases.
\end{itemize}

\section*{Treatment Options}
\begin{itemize}
  \item Treatment depends on the cause and may include inhaled medicines for diagnosed airway disease, oxygen or other respiratory support, heart treatment, or antibiotics when a bacterial infection is confirmed or strongly suspected.
  \item Do not use cough, decongestant, or inhaled medicines as a substitute for evaluation of new or worsening breathlessness.
  \item Referral to pulmonology or ENT is indicated for chronic, recurrent, or treatment-refractory cases.
\end{itemize}

\section*{Self-Care and Home Management}
\begin{itemize}
  \item Do not attempt to manage sudden or severe breathlessness at home; call local emergency services and sit upright while waiting for help.
  \item If asthma or COPD is already diagnosed, use the prescribed reliever inhaler according to the personal action plan; do not use someone else’s inhaler.
  \item Avoid smoke and other known triggers, but do not let home measures delay assessment of new or worsening symptoms.
\end{itemize}

\section*{References}
\begin{thebibliography}{99}
\bibitem{shortness_of_breath:uptodate-dyspnea} Schwartzstein RM, et al. Approach to the patient with dyspnea. In: UpToDate, Post TW (Ed), UpToDate, Waltham, MA; 2025.
\bibitem{shortness_of_breath:nejm-dyspnea} Parshall MB, et al. ``Dyspnea.'' \emph{N Engl J Med}. 2023;389(5):462-472.
\bibitem{shortness_of_breath:bmj-dyspnea} Berliner D, et al. ``Dyspnea in adults.'' \emph{BMJ}. 2024;386:e081543.
\bibitem{shortness_of_breath:harrison-dyspnea} Longo DL, et al. \emph{Harrison\'s Principles of Internal Medicine}. 21st ed. McGraw-Hill; 2022. Chapter 36: Dyspnea.
\bibitem{shortness_of_breath:cochrane-dyspnea} Bausewein C, et al. ``Non-pharmacological interventions for breathlessness.'' \emph{Cochrane Database Syst Rev}. 2023;(8):CD013987.
\bibitem{shortness_of_breath:chest-dyspnea} Mahler DA, et al. ``Mechanisms and measurement of dyspnea.'' \emph{Chest}. 2024;165(4):889-901.
\end{thebibliography}
